\section{Contextual causal language and the intervention a bank will buy}
\label{exec23:context-scope}

A causal sentence in a service note is useful evidence about what was said,
but the intervention a bank will buy requires an additional scientific
argument. Zhao, Yan and Hu's \emph{Contextual Causality with Large Language
Models: A Survey}, inspected in its 18 September 2026 arXiv version, separates
semantic, intervention and counterfactual tasks within textual context. Its
page 3 explicitly distinguishes intervening in a textual description from
manipulating the underlying data-generating process. The paper's semantic
``identification'' of causal relations should therefore be read as a language
task, rather than statistical identification of an effect. The authors state
that their text-only taxonomy organizes causal concepts and does not extend
causal theory.\citep{exec23_contextual_causality_survey}

For banking translation, the useful outputs are candidate entities, proposed
temporal relations, alternative explanations and the passages supporting each
proposal. An analyst may translate ``customers stayed after we called'' into
an action, outcome and eligible population. The same note may omit that
relationship managers called customers they already expected to retain. A
language model can make both possibilities explicit. A fluent restatement,
or a correct answer to a question about the note, cannot determine which
mechanism generated the observations.

\subsection{The same observed relationship can reverse the investment decision}
\label{exec23:context-twins}

An exact authored pair of causal laws isolates the missing information. Let
$U\sim\operatorname{Bernoulli}(1/2)$ be an unobserved account state, and let the
natural binary contact rule be $A=U$. Let $Y(a)$ be gross monetary response
under assigned contact $a$. Consider
\begin{equation}
 \mathcal M_1:\ Y(a)=100a,
 \qquad
 \mathcal M_2:\ Y(a)=100U.
 \label{exec23:context-laws}
\end{equation}
Both laws generate exactly
\begin{equation}
 P(A=0,Y=0)=\tfrac12,
 \qquad P(A=1,Y=100)=\tfrac12.
 \label{exec23:context-observed}
\end{equation}
In $\mathcal M_1$, contact causes the gross response. In $\mathcal M_2$ the
latent state causes both natural contact and response, and contact has no
effect. Under an intervention the distribution of $U$ remains unchanged.
Hence the average gross effect is USD 100 in the first law and zero in the
second.\citep{exec23_context_reference}

At USD 10 per contact, changing from no contact to contacting every eligible
account has incremental net value
\begin{equation}
 E\{Y(1)-Y(0)\}-10=
 \begin{cases}
 90,&\mathcal M_1,\\
 -10,&\mathcal M_2.
 \end{cases}
 \label{exec23:context-net}
\end{equation}
All values and mechanisms here are stipulated; they are not bank observations
or outputs of an LLM. The target compares two fixed policies, all-contact and
no-contact, rather than comparing either policy with the natural contact rule.
The witness proves that the displayed observed law alone cannot identify this
decision. It does not prove that every useful causal question is unidentified.

An assistant should preserve both graph possibilities when the supplied
evidence cannot distinguish them. A new randomized contact assignment with
adequate support would distinguish these two gross-response laws under its
design assumptions. Observing the latent state alone does not create the
missing within-state contact variation because $A=U$ under the natural rule.
Alternatively, a justified exclusion of the latent common-cause explanation
would change the model class, but that restriction needs evidence. Repeatedly
asking the same model to rewrite the same observations does not add either
source of information. The proposed business-to-graph flow should report the
unresolved distinction and the experiment or restriction that could resolve it.

\subsection{A survey can guide the evaluation design without qualifying every cited method}
\label{exec23:context-evaluation}

The survey discusses hallucinated or neglected relations, complex contexts,
multiple-choice shortcuts and limitations of external knowledge and agent
coordination. These are the survey authors' synthesis of their cited work;
those originals require separate inspection before their empirical claims
can support a package decision. Table 3 draws many representative benchmarks
from news, Wikipedia, stories, human questions and agent traces, and includes
FinCausal 2020. The presence of financial text is not evidence of contact-policy
identification or monetary utility in banking. Pages 8--9 call for extrinsic
decision quality, finer evaluation and longer agent horizons. This is a
research agenda, not an authenticated exhaustive absence result.
\citep{exec23_contextual_causality_survey}

A public or generated prototype can evaluate the resulting proposal before
bank access. Supply notes with controlled omissions, negation, reversed
chronology and conflicting explanations; separately score relation extraction,
clock preservation, graph alternatives, unsupported restrictions and appropriate
refusal. Pair identical observations from the two laws above and check that an
assistant does not invent information to force a unique effect. When an
assignment design or restriction is supplied, evaluate the changed answer and
its stated dependence on that premise. These proposed tests assess different
tasks from factual question answering. Their known-law monetary scoring is
conditional on the authored mechanisms and cannot establish actual bank value.
